\documentclass[secthm,seceqn,amsthm,ussrhead,reqno,12pt]{amsart}
\usepackage[utf8]{inputenc}
\usepackage[english]{babel}
\usepackage[symbol]{footmisc}
\usepackage{amssymb,amsmath,amsthm,amsfonts,xcolor,enumerate,hyperref, comment,longtable, cleveref}

\usepackage{times}
\usepackage{cite}
\usepackage{pdflscape}
\usepackage{ulem}
\usepackage[mathcal]{euscript}
\usepackage{tikz}
\usepackage{hyperref}
\usepackage{cancel}
\usepackage{stmaryrd}

 
\usepackage{amsfonts}
\usepackage{amssymb}
\usepackage{times}
\usepackage{xcolor}
\usepackage{tkz-graph}
\usepackage{url}
\usepackage{float}
\usepackage{tasks}
%\usepackage{refcheck}

%\DeclareMathAlphabet{\mathcal}{OMS}{cmsy}{m}{n}


%%%%%%%%%%%%%%%%%%%%%%%%%%%%%%
% Paquetes extra

\usepackage{cite}
\usepackage{hyperref}



%%%%%%%%%%%%%%%%%%%%%%%%%%%%%%%%%%%%%%%%%%%%%%%%%%%%%%%%%%%%%%%%%%%%%%%%%%%%%%%%%%%%%%%%%%%%%%%% %%%%%%%%%%%%%%%%%%%%%%%%%%%%%%%%%%%%%%%%%%%%%%%%%%%%%%%%%%%%%%%%%%%%%%%%%%%%%%%%%%%%%%%%%%%%%%%%%%%%%%%%%%%%%%%%%%%%%%%%%%%%%%%%%%%%%%%%%%%%%%%%%%%%%%%%%%%%%%%%%%%%%%%%%%%%%%%
\usepackage{amsfonts}
\usepackage{amsmath}
\usepackage{eurosym}
\usepackage{geometry}


\newcommand{\red}{\color{black}}
\newcommand{\blue}{\color{blue}}
\newcommand{\black}{\color{black}}

\usepackage{caption,booktabs}

\usepackage{refcheck}


\captionsetup{
  justification = centering
}

\setcounter{MaxMatrixCols}{12}
 

\geometry{left=0.75in,right=0.75in,top=1in,bottom=1in}

\theoremstyle{plain}
\newtheorem{theorem}{Theorem}
\newtheorem{lemma}[theorem]{Lemma}
\newtheorem{proposition}[theorem]{Proposition}
\newtheorem{remark}[theorem]{Remark}
\newtheorem{example}[theorem]{Example}
\newtheorem{definition}[theorem]{Definition}
\newtheorem{corollary}[theorem]{Corollary}
\sloppy
\begin{document}
 
 
 

\noindent{\Large
A basis of   free generalized Gerstenhaber algebras}\footnote{
The work  is  supported by RSF 22-71-10001.
}
 

 \bigskip


 \bigskip

\begin{center}

 {\bf
Ivan Kaygorodov\footnote{CMA-UBI, Universidade da Beira Interior, Covilh\~{a}, Portugal; \    kaygorodov.ivan@gmail.com} \&
 Kseniia Viatkina \footnote{Saint Petersburg  University, Saint Petersburg, Russia; \ vjatkina.k@gmail.com}  }
 


 

\end{center}


\noindent {\bf Abstract:}
{\it  A basis of the free Gerstenhaber algebra generated by a linearly ordered set  is constructed.
 }

 \bigskip


 \bigskip

 
 


\section*{Introduction}
	The notion of a Poisson bracket has its origin in the works of S.D. Poisson on celestial mechanics of the beginning of the XIX century. Since then, Poisson algebras
	have appeared in several areas of mathematics, such as:  
	Poisson manifolds~\cite{L1}, 
	algebraic geometry~\cite{BLLM,GK04}, 
	noncommutative geometry \cite{ber08},
	operads,
	quantization theory~\cite{Kon03}, 
	quantum
	groups, 
	and classical and quantum mechanics.
	The study of Poisson algebras also led to other algebraic structures, such as 
	generic Poisson algebras,
	algebras of Jordan brackets and generalized Poisson algebras,
	Novikov-Poisson algebras,
	Malcev-Poisson-Jordan algebras, 
	transposed Poisson algebras,
	$n$-ary Poisson algebras,
Gerstenhaber algebras \cite{K96}, etc.
	
 
Gerstenhaber   (odd Poisson, anti-Poisson, etc)	 brackets had appeared in many physics\cite{p1,p2} and geometric \cite{p3} contexts.
One of  	the most important "pure algebraic" sources for Gerstenhaber algebras is Hochschild cohomology \cite{Ger63}.
	Recently, some new interesting results   related to the study of
	Gerstenhaber brackets on  Hochschild cohomology had appeared in \cite{W21,V16,V19,L21}.
	Subalgebras of the Gerstenhaber algebra of differential operators are described in \cite{B06}.
		The  odd Poisson algebras are coincide with Gerstenhaber algebras and more general
odd generalized Poisson superalgebras are coincide with generalized Gerstenhaber algebras.
So, the classification of complex linearly compact simple odd generalized Poisson superalgebras is given in \cite{CK10}.
The aim of the present short note is to  construct a basis of a free generalized 	Gerstenhaber algebra.
	
	\medskip 
	
	Let us define   Gerstenhaber algebras and generalized Gerstenhaber algebras. 
	\begin{definition}
	
A generalized Gerstenhaber algebra is a $\mathbb{Z}_2$-graded vector space ${\bf V}={\bf V}_0+{\bf V}_1$ with two multiplications 
$"\cdot": {\bf V}_i \otimes {\bf V}_j \to {\bf V}_{i+j}$ and 
$"\{ , \}": {\bf V}_i \otimes {\bf V}_j \to {\bf V}_{i+j-1}$. 
The degree of an element a is denoted by $|a|,$ i.e.  $|a|=i$ if and only if $a\in {\bf V}_i$. These satisfy the identities:
\begin{longtable}{lll}
	$(1)$ &  $|a\cdot b| = |a| + |b|$ & ($"\cdot"$ has degree $0$)$;$\\
	$(2)$ & $|\{a,b\}| = |a| + |b| - 1$ & ($"\{  , \}"$  has degree $-1$)$;$\\
	$(3)$ & $(a\cdot b)\cdot c = a\cdot (b\cdot c)$ & ($"\cdot"$  is associative)$;$\\
	$(4)$ & $a\cdot b = (-1)^{|a||b|}b\cdot a$ & ($"\cdot"$ is supercommutative)$;$\\
	$(5)$ & $\{a,b\} = -(-1)^{(|a|-1)(|b|-1)} \{b,a\}$ & ("odd" antisymmetry of $"\{  , \}"$ )$;$\\
	$(6)$ & $\{a,\{b,c\}\}  = \{\{ a,b \},c\} + (-1)^{(|a|-1)(|b|-1)}\{b,\{a,c\}\}$ & ("odd" Jacobi identity)$;$\\
\end{longtable}
and a Leibniz type identity 
\begin{longtable}{lll}
	$(7)$ & $\{a,b\cdot c\} = \{a,b\} \cdot c + (-1)^{(|a|-1)|b|}b\cdot \{a,c\}+(-1)^{(|a|-1)} \mathfrak {\mathfrak D}(a) \cdot b \cdot c,$ 
\end{longtable}
	with
	$\mathfrak  {\mathfrak D}(a)=\{1,a\}$.

\end{definition}
 

 
		\begin{definition}
	A generalized Gerstenhaber algebra with 	$\mathfrak  D=0$ is a  Gerstenhaber algebra.
	\end{definition}
	
	
 
	
	
 \newpage
\section*{A basis for the free Generalized Gerstenhaber algebra}
	All vector spaces are considered over a field $\mathbb{F}$
	of characteristic zero. 
	In the present note we are using a standard method for constructing of a basis of a free algebra firstly appeared in some  papers of Shirshov \cite{shirshov} and also it was used for constructing of basis of free
	Lie algebras and superalgebras \cite{chi},
	free generalized Poisson algebras \cite{K17}, etc.

\medskip

Let us now consider free unital generalized Gerstenhaber algebras.
Take the free non-associative algebra 
$D_X$
with an "even" multiplication $\cdot$
and an  "odd" multiplication $\{,\}$
over a~field~%
$\mathbb{F}$
with the set
\begin{center}
    $X=X_0 \cup X_1=\{1,x_1, \ldots, x_n, \ldots\}$
\end{center}
of even and odd generators.
Consider the ideal~%
$I$
of
$D_X$
generated by 
\begin{longtable}{ccc}
$\Big\{$  
& \multicolumn{1}{l}{$  a\cdot b-(-1)^{|a||b|}b\cdot a, \ 
(a\cdot b)\cdot c - a\cdot (b\cdot c),$}\\
&\multicolumn{1}{l}{$   1 \cdot a-a, \
\{a,b\}+(-1)^{(|a|-1)(|b|-1)} \{b,a\},  $}\\
&$
\{a,b\}\cdot c+(-1)^{(|a|-1)|b|} b\cdot \{a,c\}+(-1)^{(|a|-1)}\mathfrak {\mathfrak D}(a)\cdot b\cdot c-\{a,b\cdot c\},
$\\

&\multicolumn{1}{r}{
$\{\{a,b\},c\} +(-1)^{(|a|-1)(|b|-1)}\{b,\{a,c\}\}-\{a,\{b,c\}\}$}&$ | \ a,b,c \in D_{X_0} \cup  D_{X_1}   \   \Big\},
$
\end{longtable}
with
$\mathfrak {\mathfrak D}(a)=\{1,a\}$.
It is easy to see  that
$D_X/I$
is the free unital generalized Gerstenhaber algebra generated by $X.$

\medskip 

Denote by
${\mathfrak G}$
the free unital generalized Gerstenhaber algebra
generated by~%
$X$,
assumed to be a linearly ordered set with
$1 <x_1 < \ldots < x_n < \ldots$.
Define the set of good words as follows.
All elements of~%
$X$
are good words.
On assuming that
all good words of length less than~%
$n$
are defined,
a word~%
$w$
of length~%
$n$
in the alphabet~%
$X$
is a good word
whenever the following conditions hold:

(1)
$w=\{u,v\}$, 
where
$u$~and~$v$
are good words with
$u>v$;

(2)
if
$w=\{\{u_1,u_2\},v\}$
then
$u_2\leq v$.

Define an order on the good words of length less than or equal to~%
$n$
to be consistent with the existing order on words
and satisfy the following condition: 
if~%
$w$
is a good word and
$w=\{u,v\}$
then
$w>u$
and
$u>v$.

Define the set
$$
M=\Big\{u,\{v,v\} \ |  \ u, v 
\mbox{ are good words in the alphabet $X$ and } v \in{\mathfrak G}_0\Big\}
$$
and then a linear order on it. %the set~$M$.
Put 
$u>v$
when
$deg(u)>deg(v)$;
if
$deg(u)=deg(v)$
then for
$u=\{u_1,u_2\}$
and
$v=\{v_1,v_2\}$,
where
$u_1>u_2$
and
$v_1>v_2$,
put
$u>v$
when
$u_1>v_1$, 
and if
$u_1=v_1$
put
$u>v$
when
$u_2>v_2$.
Denote the elements of~%the set
$M$
by
$e_i$,
where 
$e_i<e_j$
for
$i<j$.

Denote by
${\mathfrak L}_X$
the free non-associative  odd superalgebra over 
$\mathbb{F}$
with a set of generators~%
$X$
and multiplication
$\{,\}_L$.
Consider the ideal~%
$J$
of %the superalgebra
${\mathfrak L}_X$
generated by the set
\begin{longtable}{llcl}
$\Big\{$ & $\{a,b\}+(-1)^{(|a|-1)(|b|-1)} \{b,a\},$\\
&$
\{\{a,b\},c\} +(-1)^{(|a|-1)(|b|-1)}\{b,\{a,c\}\}-\{a,\{b,c\}\}$ &$ | $&$ \ a,b,c \in {\mathfrak L}_{X_0} \cup  {\mathfrak L}_{X_1}   \   \Big\}.$
\end{longtable}
The quotient
${\mathfrak L}_X/J$
is a free odd Lie superalgebra (the present case is similar to free Lie superalgebras case considered in \cite{S86});
furthermore,
%As \cite{shtern} showed,
the set
$M$
is a basis for
${\mathfrak L}_X/J$.

Denote by
$K$
the ideal in
$D_X$
generated by the elements of type
$\{ x\cdot y \ |\  x,y \neq 1\}$. 
It is easy to see that
$(D_X/I)/K  \cong {\mathfrak L}_X/J$.

Put 
$$
U=\{ e_{i_1}^{k_1} \ldots e_{i_n}^{k_n} \ |\  e_{i_1} \neq 1, e_{i_m} \in M,  i_1 < \ldots < i_n \} \cup \{1\}.
$$
This set
%$U$
is closed under the multiplication
$\cdot$,
%The multiplication
%$\cdot$
which is supercommutative,
and so we can consider the elements of
$M$
as 
$e_{i_1}^{k_1} \ldots e_{i_n}^{k_n}$,
where
$k_i \neq 0$,
and if
$e_{i_j} \in{\mathfrak G}_1$
then
$k_j=1$.
Say that
an element
$e_{i_1}^{k_1} \ldots e_{i_n}^{k_n}$
is of degree
$deg_e=\sum k_i$.
For %the element
$e_i \in M$
define
$deg_L$
as the length of the word
$e_i$
in the generators~%
$X$.
The degree of an~element
$f$
is the number of generators in the word for
$f$.

The main goal of this chapter
is to find a basis for
the free unital generalized Gerstenhaber algebra 
generated by~%
$X$.
We show that
%the set~%
$U$
is a basis for the free unital generalized Gerstenhaber algebra
generated by
$\{1, x_1, \ldots, x_n, \ldots\}$.

\begin{lemma}\label{vsp}
	The set
	$U$
	is a basis of
	${\mathfrak G}$
	as a vector space.
\end{lemma}

{\bf Proof.}
Consider a degree~%
$2$
word~%
$w$.
It is easy to see that~%
$w$
is of type
$x_ix_j$
or type
$\{x_i,x_j\}$
or type
$\{1,x_i\}$,
and
$w \in U$.
We claim that
every word of degree~%
$n$
is a linear combination of the elements of~%
$U$.
Hence,
every element of the superalgebra
${\mathfrak G}$
lies in the linear span
%is a linear combination of the elements
of~%
$U$.
We justify this by induction on~%
$n$.
Assume that 
every word of degree less than~%
$n$
lies in the linear span
%is a linear combination of the elements 
of~%
$U$.
Then
$w=w_1w_2$
or
$w=\{w_1,w_2\}$,
where
$w_1$
and
$w_2$
are words of degrees less than~%
$n$, 
that is,
$w_1$
and
$w_2$
lie in the linear span
%are linear combinations of elements
of~%
$U$.
Consider two cases.

The first case:
$w_1=\sum_i \alpha_{1i} m_{1i}$
and
$w_2=\sum_i \alpha_{2i} m_{2i}$,
where
$m_{1i},m_{2i} \in U, \alpha_{1i}, \alpha_{2i} \in k$.
Hence,
$w=\sum_{i,j} \alpha_{1i}\alpha_{2j}  m_{1i}m_{2j}$,
and,
since~%
$U$
is closed under the multiplication~%
$\cdot$,
we see that
$w=\sum_i \alpha_i m_i, m_i \in U, \alpha_i \in k$.

The second case:
$w_1=\sum_i \alpha_{1i} m_{1i}$
and
$w_2=\sum_i \alpha_{2i} m_{2i},$
where
$m_{1i},m_{2i} \in U,\alpha_{1i}, \alpha_{2i} \in k$.
Therefore,
$w=\sum_{i,j} \alpha_i \alpha_j \{m_{i},m_{j}\}$.
Consider an arbitrary term
$\{m_i,m_j\}$. 
If
$m_i$
or
$m_j$
is equal to
$n_1n_2$,
where
$n_i \in U$
with
$n_i \neq 1$,
then %we observe that
$m_j$
is of type
$n_1n_2$.
Use the relation
$$
\{m_i,n_1n_2\}= \{m_i,n_1\}n_2+(-1)^{(|m_i|-1)|n_1|}n_1\{m_i,n_2\}+(-1)^{(|m_{1i}|-1)}{{\mathfrak D}(m_i)n_1n_2}.
$$
Here we observe that
the degrees of %the elements
$\{m_i,n_1\}$,
$\{m_i,n_2\}$,
and
${\mathfrak D}(m_i)$
are less than~%
$n$
and
$n_1,n_2 \in U$.
It is easy to see that
% elements
$\{m_i,n_1\}$,
$\{m_i,n_2\}$,
and
${\mathfrak D}(m_i)$
lie in the linear span
%are linear combinations of the elements
of~%
$U$;
hence,
so does~%
$w$.

If neither
$m_i$
nor
$m_j$
is of type
$n_1n_2$,
where
$n_1,n_2 \in U$
and
$n_1,n_2 \neq 1$, 
then
$m_i,m_j \in M$.
If
$m_i=m_j \in{\mathfrak G}_1$
then
$\{m_i,m_j\} \in M$.
If  
$m_i \neq m_j$
then
we can say that
$m_i>m_j$,
and if
$m_i=\{m_{1i},m_{2i}\}$
with
$m_{1i}\geq m_{2i}$
then for
$m_{2i}\leq m_j$
we have
$\{m_i,m_j\} \in M$.
For
$m_{2i}>m_j$
we have  %it follows that
\begin{equation} \label{after}
\{m_i,m_j\} = (-1)^{(|m_{2i}|+ |m_j| -1)(|m_{i1}|-1)}\{\{m_{2i}, m_j\}, m_{i1}\}- (-1)^{|m_j|(|m_{2i}|-1)}\{\{m_{i1}, m_j\}, m_{2i}\}
\end{equation}
where
$\{\{m_{1i},m_j\},m_{2i}\} \in M$.
If
$\{m_{1i},\{m_{2i},m_j\}\} \in M$
then we infer that 
$\{m_i,m_j\}$
lies in the linear span
%is a linear combination of the elements 
of~%
$U$.
Otherwise,
we can consider
$m_{1i}$
as the
$\{,\}$-product of
$m_{11i}$
and
$m_{21i}$.
Following that,
we can continue the process
using the identity (\ref{jo2}).
Observe that 
the degree of
$m_{i}$
is greater than that of
$m_{1i}$,
which
%the degree of
%$m_{1i}$
is greater than that of
$m_{11i}$,
and so on.
Therefore,
the process is finite.
\qed
%The proof of the lemma is complete.

\medskip


Given
$b=\Pi_{k=1}^n e_k^{t_k}$,
put
$\frac{b}{e_m}=e_1^{t_1} \ldots e_{m-1}^{t_{m-1}}e_m^{t_m-1} e_{m+1}^{t_{m+1}}\cdots e_n^{t_n}$.

It is easy to see that
% multline? split?
\begin{equation}
\label{ae}
\begin{aligned}
&\{a,e_k^{t_k}\}=
\{a,e_k^{t_k-1}\}e_k +(-1)^{(|a|-1)|e_k^{t_k-1}|}e_k^{t_k-1} \{a,e_k\} +(-1)^{|a|-1}{\mathfrak D}(a)e_k^{t_k}=\\
&\left(\{a,e_k^{t_k-2}\}e_k+(-1)^{(|a|-1)|e_k^{t_k-2}|}e_k^{t_k-2} \{a,e_k\}+ (-1)^{|a|-1}{\mathfrak D}(a)e_k^{t_k-1}\right)e_k+\\
& (-1)^{(|a|-1)|e_k^{t_k-1}|}e_k^{t_k-1}\{a,e_k\} +(-1)^{|a|-1}{\mathfrak D}(a)e_k^{t_k}=\dots
\end{aligned}
\end{equation}
As we can see 
\begin{equation*}
\begin{aligned}
&(-1)^{(|a|-1)|e_k^{t_k-1}|}e_k^{t_k-1} \{a,e_k\}=(-1)^{(|a|-1)|e_k^{t_k-1}|}\cdot (-1)^{|e_k^{t_k-1}|(|a|+|e_k|-1)}\{a,e_k\}e_k^{t_k-1}=\\ &(-1)^{|e_k^{t_k-1}|(2|a|+|e_k|-2)}\{a,e_k\}e_k^{t_k-1}=\\&(-1)^{|e_k^{t_k-1}||e_k|}\{a,e_k\}e_k^{t_k-1}=(-1)^{(t_k-1)|e_k|^2}\{a,e_k\}e_k^{t_k-1}=\{a,e_k\}e_k^{t_k-1} 
\end{aligned}
\end{equation*}
( $|e_k|^2=0$ because $e_k e_k = (-1)^{|e_k||e_k|} e_k e_k$). Similarly for all terms of type $(-1)^{(|a|-1)|e_k^{t_k-i}|}e_k^{t_k-i} \{a,e_k\}e_k^{i-1}$:
$$(-1)^{(|a|-1)|e_k^{t_k-i}|}e_k^{t_k-i} \{a,e_k\}e_k^{i-1}=(-1)^{|e_k^{t_k-i}|(2|a|+|e_k|-2)}=(-1)^{(t_k -i)|e_k|^2}\{a,e_k\}e_k^{t_k-1}=\{a,e_k\}e_k^{t_k-1}$$
So the expression (\ref{ae}) reduces to
\begin{equation}
\label{aec}
\{a, e_k^{t_k} \}=t_k \{a,e_k\} e_k^{t_k-1}+(-1)^{|a|-1}(t_k-1){\mathfrak D}(a)e_k^{t_k}. 
\end{equation}

Using this notation,
we obtain
$$
\{a,b\}=\{a, \Pi_{k=1}^n e_k^{t_k} \}=
\left\{a, \frac{b}{e_n^{t_n}} \right\}e_n^{t_n}+(-1)^{(|a|-1)\left|\frac{b}{e_n^{t_n}}\right|} \frac{b}{e_n^{t_n}} \left\{a, e_n^{t_n} \right\} +(-1)^{(|a|-1)}{\mathfrak D}(a)b=
$$
$$
\left\{a, \frac{b}{e_{n-1}^{t_{n-1}}e_n^{t_n}} \right\}e_{n-1}^{t_{n-1}}e_n^{t_n}+
(-1)^{(|a|-1)\left|\frac{b}{e_{n-1}^{t_{n-1}}e_n^{t_n}}\right|} \frac{b}{e_{n-1}^{t_{n-1}}e_n^{t_n}} \left\{a, e_{n-1}^{t_{n-1}} \right\} e_n^{t_n} -
$$
$$
(-1)^{(|a|-1)\left|\frac{b}{e_n^{t_n}}\right|} \frac{b}{e_n^{t_n}} \left\{a, e_n^{t_n} \right\} +2(-1)^{(|a|-1)}{\mathfrak D}(a)b=\dost
$$
$$
=\sum_{k=1}^n
(-1)^{|e_1^{t_1}\cdots e_{k-1}^{t_{k-1}}||e_k^{t_k}|}
\{a, e_k^{t_k}\} \frac{b}{e_k^{t_k}} + (-1)^{|a|-1} (n-1) {\mathfrak D}(a)b.
$$
Let's apply (\ref{aec}) for every $\{a, e_k^{t_k} \}$:
$$
\sum_{k=1}^n
(-1)^{|e_1^{t_1}\cdots e_{k-1}^{t_{k-1}}||e_k^{t_k}|}
\left( t_k\{a, e_k\} e_k^{t_k -1} + (t_k-1)(-1)^{|a|-1}{\mathfrak D}(a)e_k^{t_k} \right)\frac{b}{e_k^{t_k}} + (-1)^{|a|-1} (n-1) {\mathfrak D}(a)b.
$$
After reducing constituents, the first term takes the form
$$(-1)^{|e_1^{t_1}\cdots e_{k-1}^{t_{k-1}}||e_k^{t_k}|}
t_k\{a, e_k\} e_k^{t_k -1}\frac{b}{e_k^{t_k}}=(-1)^{|e_1^{t_1}\cdots e_{k-1}^{t_{k-1}}||e_k^{t_k}|+|e_1^{t_1}\cdots e_{k-1}^{t_{k-1}}||e_k^{t_k-1}|}t_k\{a, e_k\} \frac{b}{e_k}$$
$$=(-1)^{|e_1^{t_1}\cdots e_{k-1}^{t_{k-1}}||e_k|}t_k\{a, e_k\} \frac{b}{e_k}$$
because $|e_1^{t_1}\cdots e_{k-1}^{t_{k-1}}|(|e_k^{t_k}|+|e_k^{t_k-1}|)=|e_1^{t_1}\cdots e_{k-1}^{t_{k-1}}|(2t_k|e_k|-|e_k|)=-|e_1^{t_1}\cdots e_{k-1}^{t_{k-1}}||e_k|$.

The second term simplified like

\begin{equation*}
\begin{aligned}
&\sum_{k=1}^n
(-1)^{|e_1^{t_1}\cdots e_{k-1}^{t_{k-1}}||e^{t_k}_k|}
(t_k-1)(-1)^{|a|-1}{\mathfrak D}(a)e_k^{t_k}\frac{b}{e^{t_{k}}_k} + (-1)^{|a|-1} (n-1) {\mathfrak D}(a)b =\\
&=\sum_{k=1}^n
(-1)^{|e_1^{t_1}\cdots e_{k-1}^{t_{k-1}}||e^{t_k}_k|}\cdot (-1)^{|e_1^{t_1}\cdots e_{k-1}^{t_{k-1}}||e^{t_k}_k|}
(t_k-1)(-1)^{|a|-1}{\mathfrak D}(a)b+(-1)^{|a|-1} (n-1) {\mathfrak D}(a)b=\\
& =\sum_{k=1}^n(-1)^{|a|-1}(t_k-1){\mathfrak D}(a)b+(-1)^{|a|-1} (n-1) {\mathfrak D}(a)b= (-1)^{|a|-1} \left(\sum_{k=1}^n t_k-1\right) {\mathfrak D}(a)b
\end{aligned}
\end{equation*}

After all these simplifications, we obtain 
\begin{equation}
\{a,b\}=\sum_{k=1}^n
(-1)^{|e_1^{t_1}\cdots e_{k-1}^{t_{k-1}}||e_k|}
t_k \{a, e_k\} \frac{b}{e_k} + (-1)^{|a|-1} \left(\sum_{k=1}^n t_k-1\right) {\mathfrak D}(a)b.
\end{equation}



Define a linear mapping
$^*:U \rightarrow U$
as follows:

$(1)$
if
$w \in X$
then
$w^*=w$;

$(2)$
if
$w \in M$
then
$w^*=w$;

$(3)$
if
$a=\Pi_{k=1}^n e_k^{a_k}, b=\Pi_{k=1}^n e_k^{b_k}$
then
$$
(ab)^*= 
(-1)^{\sum_{i=1}^n |e_i^{b_i}| |e_{i+1}^{a_{i+1}}\cdots e_n^{a_n}| } \Pi_{k=1}^n e_k^{a_k+b_k}.
$$

$(4)$
Assume that
we have defined
$^*$
for all words of length less than~%
$n$.
Consider the word
$w=\{w_1,w_2\}$
with
$w_1,w_2 \in M$ 
as an element of the free Lie superalgebra
$(M, \{,\})$.
Define
$w^*$
as a linear combination of the basis elements of %the superalgebra
$(M, \{,\})$, 
that is,
$w^*$
lies in the linear span
%is a linear combination of elements 
of~%
$M$.

$(5)$
Assume that
we have defined %the mapping
$^*$
for all words of length less than~%
$m$.
Without loss of generality,
for the word
$w=\{w_1,w_2\}$
of length~%
$m$,
where
$w_1$
or
$w_2$
lies in
$U \setminus M$,
we have
$w_2 \in U \setminus M$
and
$w_2=\Pi_{k=1}^n e_k^{t_k}$.
Then
$$(\{w_1,w_2\})^*=
\sum_{k=1}^n 
(-1)^{|e_1^{t_1}\cdots e_{k-1}^{t_{k-1}}||e_k|}
t_k\left(\{w_1, e_k\}^* \frac{w_2}{e_k}\right)^*+\\
&(-1)^{|a|-1} \left(\sum_{k=1}^n t_k-1 \right) ({\mathfrak D}(w_1)^*w_2)^*,
$$

where the elements of type
$\{x,y\}^*$
in the right-hand side %of the expression
are defined above.
\begin{theorem}\label{basis}
	The set
	$U$
	is a basis of 
	${\mathfrak G}$.
\end{theorem}

{\bf Proof.}
Consider~% the set
$U$
as a vector space basis and define two new multiplications 
$*$
and
$\{,\}^*$
as
$u*v=(uv)^*$
and
$\{u,v\}^*=(\{u,v\})^*$.
This yields the new superalgebra
$(U,*,\{,\}^*)$.
We claim that
 superalgebra
$(U, *, \{,\}^*)$
it is a generalized Gerstenhaber superalgebra.
Then we infer that~%
$U$
is a basis for
${\mathfrak G}$.

It is easy to see that~%
$*$
is associative and supercommutative, 
while
$\{,\}^*$
is \textcolor{green}{superanticommutative(у нас это не так, есть для нашего соотношения слово?)}.

We claim that
$(U,*, \{,\}^*)$
satisfies the identity
\begin{eqnarray}
	\label{jo1*}
	\{a,b*c\}^*=\{a,b\}^**c+(-1)^{(|a|-1)|b|}b*\{a,c\}^*+(-1)^{(|a|-1)}{\mathfrak D}(a)^**b*c,
\end{eqnarray}
%%Taking
%%$a,b,c \in U$,
and establish it by induction on
$deg_e(a)+deg_e(b)+deg_e(c)$.

Observe that
\begin{gather*}
	\{1,  e_i *e_j \} ^* = \{1, e_i\}^**e_j+e_i* \{1,e_j\}^*,
	\\
	\{e_k, e_i *e_j \}^*= \{e_k,e_i\}^**e_j +(-1)^{(|e_k|-1)|e_i|} e_i* \{e_k, e_j\}^* +(-1)^{(|e_k|-1)}{\mathfrak D}(e_k)^** e_i*e_j
\end{gather*}
by the definition of the mapping~%
$^*$.
Hence,
if
$deg_e(a)+deg_e(b)+deg_e(c)<4$
then (\ref{jo1*}) holds.

Suppose that
$deg_e(a)+deg_e(b)+deg_e(c)>3$.
For
$b=\Pi_{k=1}^n e_k^{b_k}$
and
$c=\Pi_{k=1}^n e_k^{c_k}$
put
$$
\gamma(b,c)= \sum_{i=1}^n |e_i^{c_i}| |e_{i+1}^{b_{i+1}} \ldots e_{n}^{b_n}|
\mbox{ and }
\mu(b,c,k)=|e_1^{b_1+c_1}\cdots e_{k-1}^{b_{k-1}+c_{k-1}}||e_k|.
$$
Hence,
\begin{eqnarray}
	\label{jordok0}
	{\mathfrak D}(a)*b*c= (-1)^{\gamma(b,c)} {\mathfrak D}(a)*\Pi_{k=1}^n e_k^{b_k+c_k}.
\end{eqnarray}

We can see that
%% ?
\begin{eqnarray}
	\label{jordok1}
	(-1)^{\gamma(b,c)}\{a,b*c\}^* =
\end{eqnarray}
\begin{eqnarray*}
	\sum_{k=1}^n 
	(-1)^{\mu(b,c,k)}
	(b_k+c_k) \{a, e_k\}^** \frac{ \Pi_{m=1}^n e_m^{b_m+c_m}}{e_k}+ 
	%\\
	(-1)^{|a|-1} \left(\sum_{k=1}^n (b_k+c_k)-1\right) {\mathfrak D}(a)* \Pi_{m=1}^n e_m^{b_m+c_m}.
\end{eqnarray*}
Furthermore,
$$
\{a,b\}^**c =
\sum_{k=1}^n  
(-1)^{|e_1^{b_1}\cdots e_{k-1}^{b_{k-1}}||e_k|}
b_k\{a, e_k\}^**\frac{b}{e_k} *c +(-1)^{|a|-1} \left(\sum_{k=1}^n b_k-1\right) {\mathfrak D}(a) *  b * c=
$$
$$
\sum_{k=1}^n  
(-1)^{|e_1^{b_1}\cdots e_{k-1}^{b_{k-1}}||e_k|
	+ \sum_{i=1}^n |e_i^{c_i}| |e_{i+1}^{b_{i+1}}\cdots e_n^{b_n}|
	-|e_k||e_{1}^{c_{1}} \cdots e_{k-1}^{c_{k-1} }| }
b_k\{a, e_k\}^**\frac{\Pi_{m=1}^n e_m^{b_m+c_m}}{e_k}
$$
$$
+ (-1)^{\sum_{i=1}^n |e_i^{c_i}| |e_{i+1}^{b_{i+1}}\cdots e_n^{b_n}|} (-1)^{|a|-1} \left(\sum_{k=1}^n b_k-1\right) {\mathfrak D}(a) *  \Pi_{m=1}^n e_m^{b_m+c_m}.
$$
Therefore,
\begin{eqnarray}
	\label{jordok2}
	\{a,b\}^**c = 
\end{eqnarray}
$$
(-1)^{\gamma(b,c)} 
\left(\sum_{k=1}^n 
(-1)^{\mu(b,c,k)}
b_k \{a, e_k\}^** \frac{ \Pi_{m=1}^n e_m^{b_m+c_m}}{e_k} +(-1)^{|a|-1}  \left(\sum_{k=1}^n b_k-1\right){\mathfrak D}(a)^**  \Pi_{m=1}^n e_m^{b_m+c_m}\right).
$$

Observe that
$$
(-1)^{|c||b|}\{a,c\}^**b =
(-1)^{|c||b|}\left(\sum_{k=1}^n 
(-1)^{|e_1^{c_1}\cdots e_{k-1}^{c_{k-1}}||e_k|}
c_k\{a, e_k\}^** \frac{c}{e_k}*b + (-1)^{|a|-1} \left(\sum_{k=1}^n c_k-1\right){\mathfrak D}(a)^* *b*c\right)=
$$
$$
\sum_{k=1}^n  
(-1)^{ {\color{blue}|e_k^{c_k}||b|}+|e_k||e_1^{c_1} \ldots e_{k-1}^{c_{k-1}}| +  \sum_{i=1}^n |e_i^{b_i}| |e_{i+1}^{c_{i+1}}\cdots e_n^{c_n}|  - |e_k||e_{1}^{b_{1}} \ldots e_{k-1}^{b_{k-1}}|}  c_k\{a, e_k\}^**\frac{\Pi_{m=1}^n e_m^{b_m+c_m}}{e_k}
$$
$$
+ (-1)^{ \sum_{i=1}^n |e_i^{c_i}| |e_{i+1}^{b_{i+1}}\cdots e_n^{b_n}|}(-1)^{|a|-1} \left(\sum_{k=1}^n c_k-1\right) {\mathfrak D}(a)^* *  \Pi_{m=1}^n e_m^{b_m+c_m}.
$$
Therefore,
\begin{eqnarray}
	\label{jordok3}
	(-1)^{|c||b|}\{a,c\}^**b =
\end{eqnarray}
$$
(-1)^{\gamma(b,c)} 
\left(\sum_{k=1}^n 
(-1)^{\mu(b,c,k)}
c_k \{a, e_k\}^** \frac{ \Pi_{k=1}^n e_m^{b_m+c_m}}{e_k} +(-1)^{|a|-1}  \left(\sum_{k=1}^n c_k-1\right){\mathfrak D}(a)^**  \Pi_{m=1}^n e_m^{b_m+c_m}\right).
$$

Finally, 
adding (\ref{jordok0}) to (\ref{jordok1}),
then subtracting (\ref{jordok2}) and (\ref{jordok3}),
we see that
$(U,*,\{,\}^*)$
satisfies (\ref{jo1*}).

Now we claim that
this superalgebra satisfies the identity
\begin{eqnarray}
	\label{jo2*}
	\{a,\{b,c\}^*\}^*=\{\{a,b\}^*,c\}^*+(-1)^{|a||b|}\{b,\{a,c\}^*\}^*.
\end{eqnarray}

For
$a,b,c \in M$
%this %(\ref{jo2*})
%holds
%because
%$(M, \{,\}^*)$
%is a Lie superalgebra.
%
%Without loss of generality,
%take
%$a=a_1*a_2$
%with
%$a_1,a_2 \in U$
%and
%$a_1,a_2 \neq 1$.
%Using (\ref{jo1*}),
%we obtain 
%% ?
%$$
%\{\{a_1*a_2,b\}^*,c\}^* +(-1)^{|a||b|} \{b,\{a_1*a_2,c\}^*\}^*=
%$$
%$$
%-(-1)^{|b||a|}\{\{b,a_1\}^**a_2+(-1)^{|b||a_1|}a_1 *\{b,a_2\}^* -{\mathfrak D}(b)^*a,c\}^*-
%$$
%$$
%(-1)^{|a|(|b|+|c|)} \{b, \{c,a_1\}^**a_2 +(-1)^{|c||a_1|}a_1*\{c,a_2\}^*-{\mathfrak D}(c)^**a_1*a_2\}^*=
%$$
%$$
%(-1)^{|b||a|+|c|(|a|+|b|)}(\{c, \{b,a_1\}^*\}^**a_2+(-1)^{(|b|+|a_1|)|c|}\{b,a_1\}^**\{c,a_2\}^*-{\mathfrak D}(c)^**\{b,a_1\}^**a_2)+
%$$
%$$
%(-1)^{|b||a|+|c|(|b|+|a|)+|b||a_1|}(\{c,a_1\}^**\{b,a_2\}^*+(-1)^{|c||a_1|}a_1*\{c,\{b,a_2\}^*\}^*-{\mathfrak D}(c)^**a_1*\{b,a_2\}^*)+
%$$
%$$
%(-1)^{|b||a|+|c|(|a|+|b|)}(\{c,{\mathfrak D}(b)^*\}^**a+(-1)^{|c||b|}{\mathfrak D}(b)^**\{c,a\}^*-{\mathfrak D}(c)^**{\mathfrak D}(b)^**a) -
%$$
%$$
%(-1)^{|a|(|b|+|c|)}(\{b,\{c,a_1\}^*\}^**a_2 +(-1)^{|b|(|c|+|a_1|)}\{c,a_1\}^**\{b,a_2\}^*-{\mathfrak D}(b)^**\{c,a_1\}^**a_2)-
%$$
%$$
%(-1)^{|a|(|b|+|c|)+|c||a|}(\{b,a_1\}^**\{c,a_2\}^*+(-1)^{|b||a_1|}a_1*\{b,\{c,a_2\}^*\}^*-{\mathfrak D}(b)^**a_1*\{c,a_2\}^*)-
%$$
%$$
%(-1)^{|a|(|b|+|c|)} (\{b,{\mathfrak D}(c)^*\}^**a+(-1)^{|b||c|}{\mathfrak D}(c)^**\{b,a\}^*-{\mathfrak D}(b)^**{\mathfrak D}(c)^**a)=
%$$
%$$
%-(-1)^{|a|(|b|+|c|)}(\{\{b,a_1\}^*,c\}^*+(-1)^{|b||a_1|}\{b,\{c,a_1\}^*\}^*)*a_2
%$$
%$$
%+(-1)^{(|b|+|c|)|a_1|}a_1*(\{\{b,a_2\}^*,c\}^*+(-1)^{|b||a_2|}\{b,\{c,a_2\}^*\}^*)-\{{\mathfrak D}(b)^*,c\}^**a-\{b,{\mathfrak D}(c)^*\}^**a=
%$$
%$$
%-(-1)^{|a||bc|}(\{\{b,c\}^*,a_1\}^** a_2+(-1)^{|bc||a_1|}a_1*\{\{b,c\}^*,a_2\}^*-{\mathfrak D}(\{b,c\}^*)^**a=\{a_1*a_2, \{b,c\}^*\}^*.
%$$
%Therefore,
%the superalgebra satisfies (\ref{jo2*}),
%and the proof of the theorem is complete.
%
%\medskip
%
%Denote by
%$GerP(n)$
%the multi-homogeneous component of the free generalized Gerstenhaber algebra 
%${\mathfrak G}$
%on
%$n+1$
%generators
%$1,x_1, \ldots, x_n$.
%%Then we can prove the following
%
%\begin{Th}\label{dimen}
%	We have
%	$dim(GerP(n))=n\cdot n!$.
%\end{Th}
%
%{\bf Proof.}
%Consider %the vector space
%$GerP[1,x_1, \ldots, x_n]$
%as a~linear subspace of the free Gerstenhaber algebra
%$P[1,x_1, \ldots, x_n]$
%on
%$n+1$
%generators. 
%It is known that the dimension of the multi-homogeneous component
%of the free Gerstenhaber algebra with
%$n+1$
%generators is
%$(n+1)!$.
%Using the description of a basis for the free generalized Gerstenhaber algebra
%(Theorem~\ref{basis}), 
%we can find a basis for the multi-homogeneous component. 
%Therefore,
%a basis for the multi-homogeneous component
%of the free generalized Gerstenhaber algebra with generators
%$1,x_1, \ldots, x_n$
%and a basis for the multi-homogeneous component
%of the free Gerstenhaber algebra with the same generators
%%$1,x_1, \ldots, x_n$,
%differ only in the elements of type
%$1 \cdot y$,
%where~%
%$y$
%is an element of the basis for the free Gerstenhaber algebra with generators
%$x_1, \ldots, x_n$.
%Thus,
%$dim(GerP(n))= (n+1)!-n!=n\cdot n!$.
%\qed
%%
%
%


%\newpage
\begin{thebibliography}{99}




  

\bibitem{p2}
Batalin I., Bering K., 
Odd scalar curvature in anti-Poisson geometry, 
Physics Letters B, 663 (2008), 1--2, 132–-135.
 
 
 \bibitem{BLLM}
 Bell J., Launois S., León  Sánchez O.,  Moosa R., 
 Poisson algebras via model theory and
differential algebraic geometry,  
Journal of the European Mathematical Society,  19 (2017), 2019--2049.
 

  \bibitem{B06} 
Bratchikov A.,  
Subalgebras of the Gerstenhaber algebra of differential operators, 
Journal of Algebra and its Applications, 16 (2017),  1, 1750005, 6 pp.


 
\bibitem{p1}
 	 Bruce A., Grabowski J., 
 	 Odd connections on supermanifolds: existence and relation with affine connections, 
 	 Journal of Physics A, 53 (2020),  45, 455203, 24 pp. 
 	 
\bibitem{CK10} 
Cantarini N., Kac V., 
Classification of linearly compact simple rigid superalgebras, 
International Mathematics Research Notices, (2010),  17, 3341--3393. 

 \bibitem{chi} 
 Chibrikov E.,  
 A right normed basis for free Lie algebras and Lyndon--Shirshov words, 
Journal of Algebra, 302 (2006),  2, 593--612. 

\bibitem{Ger63}
Gerstenhaber M., 
The cohomology structure of an associative ring.
Annals of Mathematics,  78 (1963), 2, 267--288.

 \bibitem{GK04}
 Ginzburg V.,   Kaledin D., 
 Poisson deformations of symplectic quotient singularities, 
 Advances in Mathematics, 186 (2004), 1--57

\bibitem{K17} 
Kaygorodov I.,  
Algebras of Jordan brackets and generalized Poisson algebras,  
Linear and Multilinear Algebra, 65 (2017), 6, 1142--1157.



 \bibitem{Kon03}
Kontsevich M., 
Deformation quantization of Poisson manifolds, 
Letters in Mathematical Physics, 66 (2003), 157--216.



\bibitem{K96}
Kosmann-Schwarzbach Y., 
From Poisson algebras to Gerstenhaber algebras,
Annales de l'institut Fourier (Grenoble), 46 (1996),  5, 1243--1274.


\bibitem{L1}
 Li L.-C., 
 Classical $r$-matrices and compatible Poisson structures for Lax equations on Poisson algebras,
Communications in Mathematical Physics, 203 (1999), 573--592.

\bibitem{L21}   
Lopes S.,   Solotar A., 
Lie structure on the Hochschild cohomology of a family of subalgebras of the Weyl algebra, 
Journal of Noncommutative Geometry, 15 (2021), 4, 1373--1407.
 
 
 
\bibitem{shirshov}   
 Shirshov A. I.,  
 Selected works of A. I. Shirshov. Translated from the Russian by Murray Bremner and Mikhail V. Kotchetov. Edited by Leonid A. Bokut, Victor Latyshev, Ivan Shestakov and Efim Zelmanov. Contemporary Mathematicians. Birkhäuser Verlag, Basel, 2009. viii+242 pp.
 
 
\bibitem{S86} Shtern A., 
Free Lie superalgebras,
Siberian Mathematical Journal, 27 (1986),  1, 136--140.


 \bibitem{ber08}
Van den Bergh M., 
Double Poisson algebras, 
Transactions of the American Mathematical Society, 360  (2008), 11, 5711--5769.


\bibitem{V19} 
 Volkov Yu.,  
 Gerstenhaber bracket on the Hochschild cohomology via an arbitrary resolution, 
 Proceedings of the Edinburgh Mathematical Society (2), 62 (2019),  3, 817--836
 
 \bibitem{V16}
 Volkov Yu.,  
 BV differential on Hochschild cohomology of Frobenius algebras,  
 Journal of Pure and Applied Algebra, 220 (2016),  10, 3384--3402. 

\bibitem{W21}
Wang Z., 
Gerstenhaber algebra and Deligne's conjecture on the Tate-Hochschild cohomology, 
Transactions of the American Mathematical Society, 374 (2021),  7, 4537--4577.

\bibitem{p3}	
Xu P., 
Gerstenhaber algebras and BV-algebras in Poisson geometry,  
Communications in Mathematical Physics, 200 (1999),  3, 545--560.

 
 
  
 
\end{thebibliography}


\end{document}
